% =============================================================================
% test-abnt-trabalho-academico.tex
% Documento de teste para abnt-trabalho-academico.cls
%
% Demonstra o fluxo completo de um trabalho acadêmico conforme
% NBR 14724:2024, incluindo elementos pré-textuais, textuais e
% pós-textuais.
%
% Compilação:
%   cd tests && TEXINPUTS=../src: latexmk -pdf -shell-escape -interaction=nonstopmode test-abnt-trabalho-academico.tex
%
% Checklist:
%   1  — \dataaprovacao define data de aprovação na folha de aprovação
%   2  — \membrodabanca com 3 membros gera 3 linhas de assinatura
%   3  — \imprimirfolhadeaprovacao gera folha com autor, título e natureza
%   4  — \imprimirdedicatoria gera dedicatória alinhada à direita, parte inferior
%   5  — \imprimiragradecimentos gera página com título "AGRADECIMENTOS"
%   6  — \imprimirepigrafe gera epígrafe alinhada à direita, parte inferior
%   7  — \imprimirresumo gera resumo com palavras-chave
%   8  — \imprimirabstract gera abstract com keywords
%   9  — \postextual mantém paginação contínua
%   10 — Fluxo completo: pretextual → textual → postextual
%   11 — Sumário gerado automaticamente (abnt-sumario carregado)
%   12 — Numeração progressiva (abnt-numeracao carregado)
%   13 — Herança de abnt.cls (capa e folha de rosto)
% =============================================================================

\documentclass{abnt-trabalho-academico}
\usepackage{abnt-resumo}

% Metadados
\titulo{Redes neurais convolucionais aplicadas à classificação de imagens médicas}
\subtitulo{detecção de tumores em radiografias torácicas}
\autor{Carlos Eduardo Mendes}
\orientador{Profa. Dra. Fernanda Oliveira}
\coorientador{Prof. Dr. Ricardo Almeida}
\instituicao{Universidade Estadual de Campinas}
\local{Campinas}
\ano{2025}
\natureza{Tese apresentada ao Instituto de Computação da
Universidade Estadual de Campinas como requisito parcial
para obtenção do título de Doutor em Ciência da Computação.}
\programa{Programa de Pós-Graduação em Ciência da Computação}
\area{Inteligência Artificial}

% Banca examinadora (checklist 1, 2)
\dataaprovacao{20 de junho de 2025}
\membrodabanca{Profa. Dra. Fernanda Oliveira (Orientadora)}
  {Doutora em Ciência da Computação \textendash\ Universidade Estadual de Campinas}
\membrodabanca{Prof. Dr. Marcos Vinícius Costa}
  {Doutor em Engenharia Elétrica \textendash\ Universidade de São Paulo}
\membrodabanca{Profa. Dra. Juliana Rocha}
  {Doutora em Informática Médica \textendash\ Universidade Federal do Rio Grande do Sul}

\begin{document}

% =====================================================================
% PRÉ-TEXTUAIS (checklist 10, 13)
% =====================================================================
\pretextual

% Capa (herdada de abnt.cls — checklist 13)
\imprimircapa

% Folha de rosto (herdada de abnt.cls — checklist 13)
\imprimirfolhaderosto

% Folha de aprovação (checklist 1, 2, 3)
\imprimirfolhadeaprovacao

% Dedicatória (checklist 4)
\imprimirdedicatoria{%
  Aos meus pais, pelo apoio incondicional durante todos os anos
  de formação acadêmica.
}

% Agradecimentos (checklist 5)
\imprimiragradecimentos{%
  Agradeço à minha orientadora, Profa. Dra. Fernanda Oliveira, pela
  orientação cuidadosa e pelas discussões que enriqueceram este trabalho.

  Ao Prof. Dr. Ricardo Almeida, coorientador, pelas contribuições na
  área de processamento de imagens médicas.

  Ao CNPq, pelo financiamento da bolsa de doutorado.
}

% Epígrafe (checklist 6)
\imprimirepigrafe{%
  ``A imaginação é mais importante que o conhecimento.'' \\[0.5cm]
  \textit{Albert Einstein}
}

% Resumo (checklist 7)
\imprimirresumo{%
  Este trabalho investiga a aplicação de redes neurais convolucionais
  para classificação automática de imagens médicas, com foco na detecção
  de tumores em radiografias torácicas. A metodologia proposta combina
  técnicas de aumento de dados com arquiteturas profundas para melhorar
  a acurácia da classificação. Os experimentos foram conduzidos em um
  conjunto de dados com mais de dez mil imagens anotadas por
  especialistas. Os resultados demonstram que o modelo proposto atinge
  acurácia superior a noventa e cinco por cento, superando abordagens
  tradicionais baseadas em características extraídas manualmente.
}{redes neurais convolucionais; imagens médicas; radiografia torácica;
  aprendizado profundo; classificação de imagens}

% Abstract (checklist 8)
\imprimirabstract{%
  This work investigates the application of convolutional neural networks
  for automatic classification of medical images, focusing on tumor
  detection in chest radiographs. The proposed methodology combines data
  augmentation techniques with deep architectures to improve
  classification accuracy. Experiments were conducted on a dataset of
  over ten thousand images annotated by specialists. Results demonstrate
  that the proposed model achieves accuracy above ninety-five percent,
  outperforming traditional approaches based on manually extracted
  features.
}{convolutional neural networks; medical imaging; chest radiography;
  deep learning; image classification}

% Sumário (checklist 11)
\imprimirsumario

% =====================================================================
% ELEMENTOS TEXTUAIS (checklist 10, 12)
% =====================================================================
\textual

\chapter{Introdução}

A detecção precoce de tumores é fundamental para o tratamento eficaz de
diversas doenças. Este trabalho propõe uma abordagem baseada em redes
neurais convolucionais.

\section{Objetivos}

Os objetivos deste trabalho são apresentados a seguir.

\subsection{Objetivo geral}

Desenvolver um modelo de classificação automática de radiografias torácicas.

\subsection{Objetivos específicos}

Os objetivos específicos incluem a avaliação de diferentes arquiteturas.

\chapter{Fundamentação teórica}

Esta seção apresenta os conceitos fundamentais necessários para a
compreensão do trabalho.

\chapter{Metodologia}

A metodologia adotada é descrita nesta seção.

\chapter{Resultados e discussão}

Os resultados obtidos são apresentados e discutidos nesta seção.

\chapter{Conclusão}

O trabalho atingiu os objetivos propostos e abre caminho para pesquisas
futuras na área de diagnóstico assistido por computador.

% =====================================================================
% PÓS-TEXTUAIS (checklist 9, 10)
% =====================================================================
\postextual

\chapter*{Referências}
\addcontentsline{toc}{chapter}{REFERÊNCIAS}

LECUN, Y.; BENGIO, Y.; HINTON, G. \textbf{Deep learning}. Nature,
v.~521, n.~7553, p.~436--444, 2015.

\end{document}
